\begin{preface}{Abstract}
  Most people who train do so without anyone watching them, and the faults
  that matter are the ones they cannot see themselves committing. A correction
  is only useful while the repetition is still happening, which gives a system
  about three hundred milliseconds from a fault appearing to the trainee
  hearing about it. No surveyed product delivers spoken coaching inside that
  window, publishes how its accuracy was evaluated, or states how a new
  exercise is added.

  FORMA is a mobile system that observes a trainee through the phone's camera
  and speaks a corrective cue during the set. This first phase delivers the
  perception and assessment path for three exercises: the bodyweight squat,
  the push-up and the front plank. Pose estimation runs entirely on the
  handset, movement assessment is geometric constraints declared per exercise
  rather than a learned classifier, and feedback is a pre-authored spoken cue
  selected by a discrete fault code with no language model in the live path.
  An exercise is added by authoring one specification file, which a single
  generic evaluator executes.

  The pose engine was chosen by measuring eight configurations of two engines
  on the target handset. The chosen configuration sustains 29.57 frames per
  second with a 95th-percentile per-frame latency of 27 milliseconds, against
  requirements of 22 frames per second and 65 milliseconds. Across ten
  recorded sets totalling 9,893 frames the engine returned a pose on every
  frame. Repetition counting was exact on all seven counting sets. Six
  technique faults were evaluated against the recordings. Precision was 0.990
  on the five faults judged frame by frame, no individual fault fell below
  0.957 against a requirement of 0.90, and every scripted condition was
  detected.

  Five of the six faults carry a criterion cited to published work; no
  threshold carries a source, and the specification format records that gap in
  its own fields rather than concealing it. The evaluation rests on one
  subject in one session, and an independent labelling pass by a certified
  trainer is agreed and outstanding. Both are stated as the boundary of the
  claim rather than left to be inferred.

  \vspace{1em}
  \noindent\textbf{Keywords:} on-device pose estimation, movement
  assessment, exercise technique feedback, declarative rule specification,
  mobile computer vision.
\end{preface}
